% Part: first-order-logic
% Chapter: natural-deduction
% Section: derivations

\documentclass[../../../include/open-logic-section]{subfiles}

\begin{document}

\iftag{FOL}
      {\olfileid{fol}{ntd}{der}}
      {\olfileid{pl}{ntd}{der}}

\olsection{\usetoken{P}{derivation}}

\begin{explain}
ધારણા શું છે તે આપણે કહ્યું છે અને અનુમાનના નિયમો પણ આપ્યા છે.
સ્વાભાવિક નિગમનમાં !!^{derivation}s આ બંનેમાંથી આગમનાત્મક રીતે બને છે:
દરેક !!{derivation} કાં તો એકલી ધારણા છે, અથવા એક, બે કે ત્રણ
!!{derivation}s પછી આવતું યોગ્ય અનુમાન ધરાવે છે.
\end{explain}

\begin{defn}[!!^{derivation}]
ધારણાઓ~$\Gamma$માંથી !!a{sentence}~$!A$ની \emph{!!{derivation}} એ
!!{sentence}sનું પરિમિત વૃક્ષ છે જે નીચેની શરતો સંતોષે છે:
\begin{enumerate}
\item વૃક્ષનાં સૌથી ઉપરનાં !!{sentence}s કાં તો $\Gamma$માં હોય છે,
  અથવા વૃક્ષના કોઈ અનુમાનથી !!{discharged} થાય છે.
\item વૃક્ષનું સૌથી નીચેનું !!{sentence}~$!A$ છે.
\item વૃક્ષમાં સૌથી નીચેના વાક્ય~$!A$ સિવાયનું દરેક !!{sentence},
  એવા અનુમાન-નિયમના યોગ્ય ઉપયોગનું આધારવાક્ય છે જેનું ફલિતવાક્ય
  વૃક્ષમાં એ !!{sentence}ની તરત નીચે આવે છે.
\end{enumerate}
પછી આપણે $!A$ને !!{derivation}નું \emph{ફલિતવાક્ય} અને $\Gamma$ને
તેની !!{undischarged} ધારણાઓ કહીએ છીએ.

$\Gamma$માંથી $!A$ની કોઈ !!a{derivation} અસ્તિત્વમાં હોય, તો આપણે
કહીએ કે $!A$ એ~$\Gamma$માંથી \emph{!!{derivable}} છે; પ્રતીકોમાં:
$\Gamma \Proves !A$. જો~$!A$ની એવી !!a{derivation} હોય જેમાં દરેક
ધારણા !!{discharged} હોય, તો આપણે~$\Proves !A$ લખીએ છીએ.
\end{defn}

\begin{ex}
દરેક એકલી ધારણા પોતે !!a{derivation} છે. તેથી, દાખલા તરીકે, એકલું
$!A$ !!a{derivation} છે અને એકલું $!B$ પણ છે. માનો કે આ બે પર
$\Intro{\land}$ નિયમ લાગુ પાડીએ; તો તેમાંથી નવી !!{derivation} મળે:
\begin{prooftree}
\AxiomC{$!A$}
\AxiomC{$!B$}
\RightLabel{\Intro{\land}}
\BinaryInfC{$!A \land !B$}
\end{prooftree}
આ નિયમો સામાન્ય છે: તેમાંનાં $!A$ અને~$!B$ને કોઈ પણ !!{sentence}sથી,
દાખલા તરીકે $!C$ અને~$!D$થી, બદલી શકીએ છીએ. ત્યારે ફલિતવાક્ય
$!C \land !D$ થશે; તેથી
\begin{prooftree}
\AxiomC{$!C$}
\AxiomC{$!D$}
\RightLabel{\Intro{\land}}
\BinaryInfC{$!C \land !D$}
\end{prooftree}
યોગ્ય !!{derivation} છે. અલબત્ત, ધારણાઓની અદલાબદલી પણ કરી શકીએ,
જેથી $!D$ એ~$!A$ની અને $!C$ એ~$!B$ની ભૂમિકા ભજવે. તેથી
\begin{prooftree}
\AxiomC{$!D$}
\AxiomC{$!C$}
\RightLabel{\Intro{\land}}
\BinaryInfC{$!D \land !C$}
\end{prooftree}
પણ યોગ્ય !!{derivation} છે.

હવે આપણે બીજો નિયમ, માનો કે $\Intro{\lif}$, લાગુ પાડી શકીએ છીએ. તે
આપણને પ્રેરણનું ફલિત કાઢવાની અને એ પ્રેરણના પૂર્વવર્તી સાથે તાદાત્મ્ય
ધરાવતી કોઈ પણ ધારણાને !!{discharge} કરવાની છૂટ આપે છે. તેથી નીચેનાં
બંને યોગ્ય !!{derivation}s છે:
\begin{prooftree}
\AxiomC{$\Discharge{!C}{1}$}
\AxiomC{$!D$}
\RightLabel{\Intro{\land}}
\BinaryInfC{$!C \land !D$}
\DischargeRule{\Intro{\lif}}{1}
\UnaryInfC{$!C \lif (!C \land !D)$}
\DisplayProof\bottomAlignProof
\AxiomC{$!C$}
\AxiomC{$\Discharge{!D}{1}$}
\RightLabel{\Intro{\land}}
\BinaryInfC{$!C \land !D$}
\DischargeRule{\Intro{\lif}}{1}
\UnaryInfC{$!D \lif (!C \land !D)$}
\end{prooftree}
તેઓ અનુક્રમે બતાવે છે કે $!D \Proves !C \lif (!C \land !D)$ અને
$!C \Proves !D \lif (!C \land !D)$.

યાદ રાખો કે ધારણાઓ નિવૃત્ત કરવી એ છૂટ છે, આવશ્યકતા નહીં: આપણે ધારણાઓ
નિવૃત્ત કરવી જ પડે એમ નથી. ખાસ કરીને, ધારણાઓ !!{derivation}માં હાજર ન
હોય તોપણ નિયમ લાગુ કરી શકીએ છીએ. દાખલા તરીકે, !!{discharged} કરવા માટે કોઈ
ધારણા~$!A$ ન હોવા છતાં નીચેનું માન્ય છે:
\begin{prooftree}
  \AxiomC{$!B$}
  \DischargeRule{\Intro{\lif}}{1}
  \UnaryInfC{$!A \lif !B$}
\end{prooftree}
\end{ex}

\end{document}
